\documentclass[11pt,a4paper]{article}
\usepackage[spanish,es-noshorthands]{babel}
\usepackage[utf8]{inputenc}
\usepackage[T1]{fontenc}
\usepackage{lmodern}
\usepackage{geometry}
\geometry{margin=2.5cm}
\usepackage{hyperref}
\usepackage{booktabs}
\usepackage{longtable}
\usepackage{verbatim}
\usepackage{enumitem}

\hypersetup{colorlinks=true,linkcolor=blue,urlcolor=blue}
\setlength{\parskip}{0.6em}
\setlength{\parindent}{0pt}

\title{Análisis léxico del lenguaje Prolog:\\investigación, diseño,
implementación y validación\\{\large INFO1148, Teoría de la Computación}}
\author{[PENDIENTE: completar con los nombres reales antes de entregar]}
\date{17 de septiembre de 2026 --- Prof.\ M.\ Lévano}

\begin{document}
\maketitle

\begin{center}
  \begin{tabular}{ll}
    \toprule
    Campo & Valor \\
    \midrule
    Curso & INFO1148, Teoría de la Computación \\
    Repositorio & \url{https://github.com/FernandoValdes01/Tarea01TC} \\
    \bottomrule
  \end{tabular}
\end{center}

\section*{Resumen ejecutivo}
Este informe presenta el análisis, diseño, implementación y
validación de un analizador léxico para un subconjunto de Prolog
inspirado en la notación de ISO Prolog y SWI-Prolog. El trabajo
define formalmente el catálogo de tokens con una expresión regular
por categoría, modela las categorías principales con autómatas
finitos, muestra la determinización por construcción de subconjuntos
y la minimización por refinamiento de particiones sobre un
subconjunto representativo, e implementa el lexer en Python con
recorrido de izquierda a derecha, máxima coincidencia y prioridades
explícitas. El lexer emite tipo, lexema original, línea y columna por
token, administra una tabla de lexemas con deduplicación por par tipo
y lexema, y reporta seis clases de error con fragmento y posición
continuando el análisis. La validación usa una suite de 85 pruebas
automatizadas con 26 casos válidos, 22 de prioridad, 16 inválidos y
dos archivos completos, con resultado real de 85 aprobadas y cero fallos.

\textbf{Palabras clave:} análisis léxico, Prolog, expresiones
regulares, autómatas finitos, tabla de lexemas.

\section{Introducción}
El análisis léxico es la primera fase de un compilador: transforma la
cadena fuente en una secuencia ordenada de tokens con posición,
separando las decisiones regulares (qué formas son válidas) de las
decisiones sintácticas (en qué orden pueden aparecer). Un lexer
correcto y bien diagnosticado simplifica el parser, produce errores
comprensibles y deja evidencia reproducible de cada decisión.

Este informe documenta el subconjunto de Prolog reconocido, su
especificación formal, los autómatas que lo modelan, la
implementación en Python, la tabla de lexemas, el tratamiento de
errores y el plan de pruebas con resultados reales. Cada afirmación
remite a su respaldo: sección de la especificación, módulo del
código, prueba automatizada o salida de un comando ejecutado.

\section{Objetivos}
El objetivo general es analizar, diseñar, implementar y validar un
analizador léxico para un subconjunto definido de Prolog, aplicando
alfabetos, lenguajes regulares, expresiones regulares y autómatas
finitos, con pruebas reproducibles y diagnósticos verificables.

Los objetivos específicos medibles son: definir formalmente cada
categoría con nombre de token, patrón, ejemplos válidos e inválidos y
atributo asociado; construir una expresión regular por categoría
resolviendo superposiciones con máxima coincidencia y prioridades;
modelar las categorías principales con AFN o AFD y presentar un AFD
integrado o una estrategia equivalente con estados de aceptación
identificados; mostrar la determinización por construcción de
subconjuntos y la minimización con justificación de equivalencias
sobre un subconjunto representativo; implementar el lexer con
posiciones 1-based, tabla de lexemas y seis clases de error
recuperables; y validar con al menos 20 pruebas válidas, 8 inválidas,
casos de prioridad y dos archivos completos, registrando el resultado
real de la suite.

\section{Alcance del analizador}
El analizador reconoce exactamente el catálogo de
\verb|docs/informe/01_especificacion_lexica.md|: átomos simples y
entre comillas simples, variables nombradas y anónima, enteros y
reales sin signo, strings con comillas dobles, los operadores de las
seis familias (cláusula, consulta, DCG, unificación y comparación,
aritméticos, control), los nueve delimitadores de un carácter y los
elementos ignorables (espacios, tabulaciones, saltos de línea y
comentarios de línea y de bloque).

Quedan explícitamente fuera del alcance el análisis sintáctico y
semántico: no se verifica el orden de los tokens, no se construye un
árbol sintáctico, no se comprueba el balance de paréntesis, corchetes
o llaves, no se resuelven ámbitos ni variables repetidas, no se
implementa unificación, no se evalúan expresiones y no se decide si
una secuencia de tokens forma un programa válido. Una \verb|)|
aislada es un token \verb|RPAREN| correcto; su posible improcedencia
pertenece a una fase posterior.

\section{Especificación léxica}
El contrato completo vive en
\verb|docs/informe/01_especificacion_lexica.md| y este capítulo lo
resume sin modificarlo. El lexema de cada token es el segmento
original de entrada, sin normalización ni conversión de números y sin
des-escapar literales. El alfabeto de identificadores es ASCII y
cualquier letra, dígito o espacio Unicode fuera de un literal es un
carácter no admitido.

{\small
\begin{longtable}{llll}
  \toprule
  Token & Patrón Python (\verb|re|) & Válidos & Exclusiones \\
  \midrule
  \endhead
  \verb|ATOM| & \verb|[a-z][a-z0-9_]*| & \verb|padre|,
  \verb|persona_1|, \verb|isla| & \verb|Padre|, \verb|á| \\
  \verb|QUOTED_ATOM| & \verb|'...\\['nrt]...| & \verb|'Juan Pérez'|,
  \verb|':-'|, \verb|'it\'s'| & sin cierre, \verb|'a\q'| \\
  \verb|VARIABLE| & \verb|[A-Z]...| o \verb|_[...]+| & \verb|X|,
  \verb|Persona|, \verb|_Temporal| & \verb|_| solo, \verb|9X| \\
  \verb|ANONYMOUS| & \verb|_| & \verb|_| & \verb|_X|, \verb|__| (son
  VARIABLE) \\
  \verb|INTEGER| & \verb|[0-9]+| & \verb|0|, \verb|25|, \verb|0007| &
  \verb|5.|, \verb|.5| \\
  \verb|REAL| & \verb|[0-9]+\.[0-9]+| & \verb|3.14|, \verb|0.5| &
  \verb|5.|, \verb|1.2.3|, exponentes \\
  \verb|STRING| & \verb|"...\\["nrt]...| & \verb|"hola"|,
  \verb|"línea\n"| & sin cierre, \verb|"a\q"| \\
  \verb|OPERATOR| & unión normativa (Sec.\ 8) & \verb|:-|,
  \verb|?- |, \verb|-->|, \verb|\==|, \verb|=..| & \verb|:|, \verb|?|,
  \verb|\| aisladas \\
  Delimitadores & un carácter & \verb|()[]{}|, \verb|,| & sin
  validación de balance \\
  \bottomrule
\end{longtable}
}

Las familias de \verb|OPERATOR| son \verb|CLAUSE| (\verb|:-|),
\verb|QUERY| (\verb|?-|), \verb|DCG| (\verb|-->|),
\verb|UNIFICATION_COMPARISON| (\verb|=|, \verb|\=|, \verb|==|,
\verb|\==|, \verb|=..|, \verb|<|, \verb|=<|, \verb|>|, \verb|>=|),
\verb|ARITHMETIC| (\verb|+|, \verb|-|, \verb|*|, \verb|/|, \verb|//|,
\verb|**|, \verb|is|, \verb|mod|) y \verb|CONTROL| (\verb|\+|,
\verb|!|, \verb|;|). La unión normativa implementada en
\verb|src/lexer.py| (\verb|OPERATOR_RE|) es idéntica a la del
contrato y ordena cada prefijo largo antes que su prefijo corto, con
guardas de palabra \verb|(?![A-Za-z0-9_])| para \verb|is| y \verb|mod|.

Las convenciones adoptadas son: los números no tienen signo
(\verb|-12| son dos tokens); \verb|INTEGER| y \verb|REAL| no admiten
exponentes ni puntos sin dígitos a ambos lados; los escapes válidos
son exactamente barra inversa más uno permitido según el delimitador
(\verb|\\|, \verb|\'|, \verb|\"|, \verb|\n|, \verb|\r|, \verb|\t|), y
cualquier otra pareja es \verb|INVALID_ESCAPE|; ningún literal admite
un salto de línea real en su interior; \verb|%| inicia el comentario
% de línea hasta antes del salto; \verb|/*| abre un bloque que
% termina en el primer \verb|*/| sin anidamiento; y las posiciones
% son 1-based con tabulación de una columna y \verb|\r\n| contado
% como un único salto.

\section{Diseño de autómatas}
Las figuras usan la convención \verb|estado --símbolo--> estado|, con
\verb|(( ))| para aceptación y la etiqueta del token reconocido.
Abreviaturas: \verb|LOWER=[a-z]|, \verb|UPPER=[A-Z]|,
\verb|DIGIT=[0-9]|, \verb|ID=[A-Za-z0-9_]|.

\begin{verbatim}
Figura 1 (átomos): q0 --LOWER--> ((q1 ATOM)); q1 --[a-z0-9_]--> q1.
Figura 2 (variables): q0 --UPPER--> ((q1 VARIABLE)); q0 --_--> q2;
  q1 --ID--> q1; q2 --ID--> ((q3 VARIABLE)); q3 --ID--> q3;
  q2 sin continuación = ANONYMOUS_VARIABLE por prioridad.
Figura 3 (números): q0 --DIGIT--> ((q1 INTEGER)); q1 --DIGIT--> q1;
  q1 --.---> q2; q2 --DIGIT--> ((q3 REAL)); q3 --DIGIT--> q3;
  q0 --.---> q4; q4 --DIGIT--> q5; q5 --[0-9.]--> q5 (trampa).
Figura 4 (operadores): s0 con ramas \ -> {\=, \==, \+},
  = -> {=, ==, =.., =<}, / -> {/, //, /*}, * -> {*, **},
  : -> {:-}, ? -> {?-}, - -> {-->}, más <, >, +, !, ;;
  cada aceptación porta lexema y familia.
Figura 5 (ignorables): [ \t] en ciclo de una columna; \n, \r, \r\n
  un salto; % hasta antes del salto; /* ... */ hasta el primer */
  o UNTERMINATED_BLOCK_COMMENT.
Figura 6 (integrado): despachador por primera clase + máxima
  coincidencia + prioridades; equivalente a la unión de Figuras 1-5.
\end{verbatim}

El estado inicial de cada autómata es \verb|q0| (o \verb|s0|); los
estados de aceptación y sus tokens se indican en cada figura. La
trampa \verb|q5| captura \verb|.5| y las corridas que el validador
reporta como \verb|MALFORMED_NUMBER|; el lexer implementa este
autómata como corrida máxima de \verb|[0-9.]+| seguida de validación
exacta contra \verb|INTEGER| o \verb|REAL|.

\section{Determinización y minimización}
Para la determinización se tomó el subconjunto representativo de
ambigüedad entre operador y comentario con las ramas \verb|/|,
\verb|//| y \verb|/* ... */|, construido por Thompson como unión de
tres ramas con transiciones vacías desde un inicio común. La
$\varepsilon$-cerradura del inicio es el conjunto de los tres primeros estados de
rama más el inicio, y como ninguna rama intermedia usa transiciones
vacías adicionales, cada $\varepsilon$-cerradura posterior es la identidad; se
declara así porque el cálculo se hizo y su resultado es trivial, no
porque se haya omitido.

\begin{center}
\begin{tabular}{lllll}
  \toprule
  Subconjunto & Con \verb|/| & Con \verb|*| & Con otro & Aceptación \\
  \midrule
  S0 (inicio) & S1 & --- & --- & ninguna \\
  S1 (tras \verb|/|) & S2 & S3 & operador \verb|/| & \verb|OPERATOR('/')| \\
  S2 (tras \verb|//|) & --- & --- & fin & \verb|OPERATOR('//')| \\
  S3 (cuerpo \verb|/*|) & S3 & S4 & S3 & ninguna \\
  S4 (tras \verb|*|) & cierre si sigue \verb|/| & S4 & S3 & cierre
  con \verb|*/| \\
  \bottomrule
\end{tabular}
\end{center}

La tabla muestra por qué el comentario debe intentarse antes que el
operador: desde S1 el símbolo \verb|*| compromete al autómata con el
bloque y abandonar esa rama violaría la máxima coincidencia. El lexer
respeta este orden al tratar \verb|/*| en la fase de ignorables antes
de probar operadores.

Para la minimización se tomó el AFD de números de la Figura 3 con la
partición inicial entre aceptadores (\verb|q1| entero, \verb|q3|
real) y no aceptadores. El refinamiento separa \verb|q1| de \verb|q3|
porque ante \verb|.| transitan a bloques distintos, lo que justifica
que entero y real nunca se fusionen y que la decisión entre ambos la
tome la longitud del lexema. Los estados trampa de número mal formado
resultan equivalentes entre sí porque ciclan sobre \verb|[0-9.]| sin
aceptar, de modo que se fusionan en una única trampa que el conductor
reporta como \verb|MALFORMED_NUMBER| con la corrida completa.

\section{Reglas de prioridad y máxima coincidencia}
En cada posición se elige el candidato válido más largo y solo en
empate se aplica la prioridad contractual.

{\small
\begin{longtable}{llll}
  \toprule
  Entrada & Candidatos & Regla & Resultado \\
  \midrule
  \endhead
  \verb|\==| & \verb|\==|, \verb|\=| & más largo & OPERATOR unificación (P01) \\
  \verb|=..| & \verb|=..|, \verb|=| & más largo & OPERATOR unificación (P02) \\
  \verb|==| & \verb|==|, \verb|=| & más largo & OPERATOR unificación (P03) \\
  \verb|:-| & \verb|:-| frente a \verb|:| & 2 caracteres & OPERATOR
  cláusula (P05) \\
  \verb|?- | & \verb|?- | frente a \verb|?| & 2 caracteres & OPERATOR
  consulta (P06) \\
  \verb|-->| & \verb|-->|, \verb|-| & más largo & OPERATOR DCG (P07) \\
  \verb|//| & \verb|//|, \verb|/| & más largo & OPERATOR aritmético (P08) \\
  \verb|**| & \verb|**|, \verb|*| & más largo & OPERATOR aritmético (P09) \\
  \verb|_| & anónima & regla del anónimo & ANONYMOUS (P10) \\
  \verb|_Tmp| & variable completa & más largo & VARIABLE (P11) \\
  \verb|isla| & ATOM(4) vs prefijo & guarda & ATOM (P12) \\
  \verb|is|, \verb|mod| & empate ATOM/OPERATOR & reservada & OPERATOR
  aritm. (P13-14) \\
  \verb|1.25| & REAL vs INTEGER & más largo & REAL (P16) \\
  \verb|X=Y| & adyacentes & avanza cursor & 3 tokens (P17) \\
  \verb|12abc| & INTEGER+ATOM & corrida \verb|[0-9.]| & 2 tokens, sin
  error (P21) \\
  \verb|5.|, \verb|.5|, \verb|1..2| & numérico inválido & detector &
  MALFORMED (I07-09) \\
  \bottomrule
\end{longtable}
}

La nota de \verb|12abc| es una divergencia declarada con el prompt de
pruebas, no con el contrato: la corrida numérica solo incluye dígitos
y puntos, de modo que \verb|12| es un entero válido seguido del átomo
\verb|abc|. El mínimo de 8 inválidas se cumple con otros 16 casos con
error real.

\section{Diseño e implementación}
La arquitectura tiene cuatro módulos con una sola dirección de
dependencias: \verb|tokens.py| define \verb|TokenType|,
\verb|OperatorFamily|, el mapa \verb|OPERATOR_FAMILY| y el
\verb|Token| con posición de inicio y familia opcional;
\verb|errors.py| define \verb|ErrorKind| y \verb|LexError|;
\verb|symbol_table.py| implementa \verb|SymbolTable| con
deduplicación por par tipo y lexema; \verb|lexer.py| expone
\verb|tokenize(source)| y \verb|tokenize_file(path)| y retorna un
\verb|LexResult| con tokens en orden, errores acumulados y tabla; y
\verb|main.py| ofrece la CLI que imprime cada token como \texttt{<TIPO, lexema, línea, columna>} más la familia cuando existe.

El flujo recorre la entrada de izquierda a derecha: primero consume
ignorables en bucle, luego prueba el detector numérico, después los
literales con validación de escapes, después \verb|OPERATOR_RE|
anclado al cursor, después palabras y anónima, después delimitadores
y finalmente registra \verb|INVALID_CHARACTER|. Cada token guarda la
línea y columna de su primer carácter; la tabulación avanza una
columna y \verb|\r\n| cuenta como un único salto.

La salida real del programa válido empieza con
\verb|<ATOM, 'padre', 2, 1>|, \verb|<LPAREN, '(', 2, 6>|,
\verb|<ATOM, 'juan', 2, 7>| y totaliza 207 tokens, 0 errores y 48
entradas de tabla. Los comandos reproducibles son
\verb|python -m compileall src tests|, \verb|python -m pytest -q|,
\verb|python -m src.main tests/corpus/entradas/programa_valido.pl| y
\verb|python -m src.main tests/corpus/entradas/programa_con_errores.pl|
(en este entorno, la suite se ejecutó con \verb|uv run --with pytest|).

\section{Tabla de lexemas}
Se registran \verb|ATOM|, \verb|QUOTED_ATOM| y \verb|VARIABLE| como
identificadores, más \verb|INTEGER|, \verb|REAL| y \verb|STRING| como
literales conservables, siempre con el lexema original y sin
conversiones. La clave de deduplicación es el par tipo y lexema: dos
apariciones idénticas comparten índice y dos formas textuales
distintas nunca se fusionan. Cada token registrable porta su índice
puramente léxico. La variable anónima \verb|_| nunca se registra, y
tampoco operadores, delimitadores, espacios ni comentarios. La tabla
no guarda alcances, usos, unificaciones ni tipos semánticos.

\section{Tratamiento de errores léxicos}
Los seis errores contractuales se acumulan sin detener el recorrido y
cada diagnóstico incluye clase, fragmento, línea, columna y mensaje
útil. El programa con errores produce exactamente estos 12
diagnósticos: \verb|UNTERMINATED_QUOTED_ATOM| en 3:6,
\verb|UNTERMINATED_STRING| en 5:6, \verb|INVALID_CHARACTER('@')| en
7:6, \verb|MALFORMED_NUMBER('5.')| en 8:6,
\verb|MALFORMED_NUMBER('.5')| en 9:6, \verb|MALFORMED_NUMBER('1..2')|
en 10:6, \verb|INVALID_CHARACTER(':')| en 11:5,
\verb|INVALID_CHARACTER('?')| en 12:5, \verb|INVALID_CHARACTER| de
barra inversa en 13:5, \verb|INVALID_ESCAPE| en 14:8 y 15:8, y
\verb|UNTERMINATED_BLOCK_COMMENT| en 17:1 hasta EOF. El bloque sin
cierre no admite tokens posteriores; todos los demás errores dejan
tokens válidos después.

\section{Plan de pruebas y resultados}
La suite \verb|tests/test_lexer.py| suma 85 pruebas: 3 del modelo
\verb|Token|, 26 válidas con tokens y posiciones exactas, 22 de
prioridad, 16 inválidas con clase, fragmento, posición y
recuperación, 12 transversales de posiciones, comentarios, escapes,
tabla y acumulación, y 6 de archivos incluyendo los dos programas
completos y los corpus en disco. La matriz completa con el resultado
real de cada caso vive en \verb|docs/informe/04_matriz_de_pruebas.md|.

El archivo válido produce 207 tokens, 0 errores y 48 entradas,
cubriendo los 17 tipos de token y las 6 familias. El archivo con
errores produce 69 tokens, 12 errores de las 6 clases en 12 líneas
distintas y 28 entradas. El comando ejecutado fue
\verb|uv run --with pytest python -m pytest -q| con resultado real
\verb|85 passed in 0.08s|; además \verb|python -m compileall src tests| terminó sin errores y \verb|git diff --check| salió limpio.

\section{Análisis de resultados}
La cobertura es total sobre el catálogo: cada tipo de token y cada
familia aparece en el corpus válido, cada clase de error tiene al
menos dos variantes cuando el contrato lo permite y cada regla de
prioridad tiene su prueba dedicada. Los casos límite confirman que
CRLF cuenta como un salto, que la tabulación avanza una columna, que
los comentarios dentro de literales son contenido, que \verb|-12| son
dos tokens y que \verb|)| aislado sigue siendo token.

La principal dificultad fue la interacción entre el punto de fin de
cláusula y el detector numérico: una cláusula que termina en número
sin espacio es \verb|MALFORMED_NUMBER| por contrato, de modo que los
archivos válidos separan con espacio. La segunda fue \verb|12abc|,
resuelta como dos tokens según la definición de corrida numérica. La
recuperación quedó demostrada con 69 tokens válidos alrededor de 11
errores recuperables. El único límite real es documental:
\verb|agents.md| y la plantilla de informe no se encontraron en las
rutas indicadas.

\section{Conclusiones}
El trabajo muestra que un lexer riguroso nace de decisiones
contractuales explícitas antes que del código: alfabeto ASCII,
números sin signo, escapes cerrados por delimitador, comentarios no
anidables y máxima coincidencia con prioridades resolvieron todas las
ambigüedades sin casos especiales en la implementación. Los objetivos
se cumplieron de forma medible con 85 pruebas aprobadas, cobertura
total de categorías y familias, y dos archivos completos que
ejercitan el recorrido real.

Las limitaciones reales son que el lexer solo reconoce el subconjunto
contratado, que los archivos válidos deben separar con espacio un
número del punto final, y que el informe conserva marcadores
pendientes para integrantes y plantilla. Como mejoras futuras se
proponen modos de salida adicionales (JSON o CSV de tokens), un
visualizador de autómatas generado desde las expresiones del código y
pruebas de rendimiento sobre archivos grandes, sin ampliar el
subconjunto ni entrar en análisis sintáctico.

\section{Referencias}
Las fuentes realmente consultadas, en formato APA 7, se detallan en
\verb|docs/informe/referencias.md| y comprenden el enunciado de la
tarea, la especificación léxica contractual, la matriz de pruebas, el
código y el corpus del repositorio, y la documentación oficial de
Python usada para implementar y ejecutar la suite. No se inventaron
editoriales, autores, URL ni fechas.

\section{Anexos}
Los anexos viven en \verb|docs/informe/anexos.md| e incluyen el
enlace real al repositorio, los listados de autómatas completos, las
salidas extensas de ambos programas, la evidencia de los comandos de
verificación y la lista de comprobación contra la rúbrica con el
archivo o prueba que respalda cada dimensión.
\end{document}
